\section{Reinforcement Learning}
\label{sec:rl}

We train the fine-tuned LLM with PPO (Proximal Policy Optimization)~\citep{schulman2017ppo} and Generalized Advantage Estimation (GAE)~\citep{schulman2016gae}.

\subsection{The environment}
\textbf{Initialization.} The environment starts with a SystemVerilog design and an assertion to prove.

\textbf{Action.} The agent's action is a DAG of assertions rooted at the target, with helper assertions as descendants, together with cutpoint, abstraction, and blackboxing commands.

\textbf{Environment mechanics.} Given an action, the environment applies the cutpoint, abstraction, and blackboxing commands to the design. It then verifies the assertions in the tree in topological order; when proving an assertion, it assumes its descendants properties are true.

\textbf{Observation.} The agent observes the design and the target property. If the prover fails, the agent also observes which assertions failed and why (e.g., counterexample or timeout).

\textbf{Termination.} The environment terminates when the target assertion is proven or when a maximum number of steps is reached.

\subsection{Prover flow}
\label{sec:prover}

The environment requires a prover that parses SVA, discharges assertions
unattended at scale, and supports the cutpoint, abstraction, and blackboxing
commands of the action space. Candidate flows:

\begin{itemize}
  \item \textbf{SymbiYosys + Verific (Tabby CAD Suite)}~\cite{symbiyosys}:
    Yosys-based flow with the industry-standard Verific frontend. Full IEEE
    1800 SVA support; cutpoints and blackboxing map directly to Yosys passes;
    engines include BMC, k-induction, and IC3 via SMT solvers and ABC.
    Proprietary but free for academic use.
  \item \textbf{SymbiYosys (open-source Yosys frontend)}~\cite{symbiyosys}:
    same flow without Verific. Fully open source, but parses only a subset of
    SVA (immediate and simple concurrent assertions), which would exclude the
    temporally deep properties in our corpus.
  \item \textbf{EBMC}~\cite{ebmc}: word-level bounded model checker with
    k-induction from the CPROVER family. Native SystemVerilog frontend with
    broader SVA coverage than open-source Yosys; free and dependency-light, so
    trivial to deploy on a cluster, but weaker engine variety and no built-in
    cutpoint support.
  \item \textbf{Commercial tools (e.g., JasperGold, VC Formal)}: strongest
    engines and full SVA, but license-server seats constrain the parallel,
    unattended proof workers our RL loop requires (Section~\ref{sec:cost}).
\end{itemize}

We use SymbiYosys with the Verific frontend as the primary flow---it is the
only free option that combines full SVA parsing with unattended operation and
native cutpoint/blackbox support---and EBMC as a fallback for designs outside
the Verific flow.

\subsection{Reward design}
Each step receives a large positive reward when the target assertion is proved. Separately, every assertion in the tree incurs a penalty based on the time the engine spent checking it and on whether it was left unproven.

\subsection{Designs and assertions for the environment}
\label{sec:rl-instances}

Each instance pairs a SystemVerilog design with a target assertion that must
satisfy \textbf{ground-truth provability}---the target is
known to hold, so failures are attributable to the agent rather than to a false
property (we use only formally verified SVA, never generated assertions of
unknown status).

\textbf{Sources.}
We reuse the verified-SVA subset of the fine-tuning corpus
(Section~\ref{sec:ft-sources}): HierSVA~\cite{hiersva2026},
OpenTitan~\cite{opentitan}, riscv-formal~\cite{riscvformal}, and HWSec-UNC
verification-benchmarks~\cite{ahmad2024hwsec}.

\textbf{Splitting and leakage control.}
Instances are partitioned by design (not assertion) into training and held-out
sets, so no module seen during RL appears at test time. Designs shared with the
fine-tuning corpus stay in the training split, keeping the held-out set a clean
measure of generalization to unseen designs.

\subsection{Curriculum learning}
\label{sec:curriculum}

\textbf{Difficulty grading.}
Each pair carries a static difficulty proxy---cone-of-influence state-bit count,
hierarchy depth of the supporting module tree, and temporal depth of the
property (nesting of \texttt{\#\#}, \texttt{|->}, and unbounded
operators)---calibrated against baseline single-shot proof time and timeout rate
and binned into tiers consumed by the curriculum.

We start the LLM on simple designs and assertions with low timeouts, and gradually increase the complexity of the designs and assertions, as well as the timeouts, as the LLM improves.
